\documentclass[tikz,border=6pt]{standalone}
\usepackage{ctex}
\usepackage{amsmath}
\usetikzlibrary{calc, arrows.meta, angles, quotes}

\begin{document}
\begin{tikzpicture}[scale=1.3, thick]

% 空间向量投影示意图
% 向量 b 在向量 a 方向上的投影

% O 为公共起点
\coordinate (O) at (0,0);

% 向量 a 沿水平方向
\coordinate (Aend) at (3.8,0);

% 向量 b 斜向（与 a 夹角 θ≈50°）
\coordinate (Bend) at (2.0,2.4);

% 画向量 a（黑色）
\draw[-{Stealth[length=8pt, width=4pt]}, black, very thick] (O) -- (Aend);
\node[right, font=\normalsize] at (Aend) {$\vec{a}$};

% 画向量 b（蓝色）
\draw[-{Stealth[length=8pt, width=4pt]}, blue, very thick] (O) -- (Bend);
\node[above, font=\normalsize, blue] at (Bend) {$\vec{b}$};

% b 的终点在 a 方向上的垂足 H
% H = 投影长度 * a方向单位向量
% |b|cos θ = 2.0（水平分量即为投影）
\coordinate (H) at (2.0,0);

% 垂线（从 Bend 到 H）
\draw[gray, dashed, thin] (Bend) -- (H);

% 直角标记
\draw[gray, thin] (H) ++(0,0.13) -- ++(0.13,0) -- ++(0,-0.13);

% 投影长度标注（O 到 H 的双箭头）
\draw[{Stealth[length=5pt]}-{Stealth[length=5pt]}, red!70, thin]
  (0,-0.35) -- (2.0,-0.35)
  node[midway, below, font=\small, red!80] {投影 $=|\vec{b}|\cos\theta$};

% 投影线（加粗红色，O 到 H）
\draw[red, very thick] (O) -- (H);
\fill[red] (H) circle [radius=2.2pt];

% 夹角弧
\draw[black, thin] (O) ++(0.5,0) arc[start angle=0, end angle=50.2, radius=0.5cm];
\node[right, font=\small] at (0.53, 0.22) {$\theta$};

% 原点
\node[below left, font=\small] at (O) {$O$};

% 公式标注
\node[draw=blue!50, fill=blue!5, rounded corners=4pt,
      align=center, font=\normalsize, inner sep=8pt] at (1.8,3.5)
  {$\vec{b}$ 在 $\vec{a}$ 方向上的投影};
\node[draw=red!50, fill=red!5, rounded corners=4pt,
      align=center, font=\normalsize, inner sep=8pt] at (1.8,2.8)
  {$= |\vec{b}|\cos\theta = \dfrac{\vec{a}\cdot\vec{b}}{|\vec{a}|}$};

\end{tikzpicture}
\end{document}
